\documentclass[10pt,a4paper]{amsart}
\usepackage[margin=19mm]{geometry}
\usepackage{amsmath,amssymb,mathtools}
\usepackage[scr=rsfs,cal=euler]{mathalfa}
\usepackage{xfrac}
\usepackage[unicode,hidelinks,bookmarks=false]{hyperref}
\newcommand{\ie}{i.e.\ }
\newcommand{\cf}{cf.\ }
\newcommand{\resp}{respectively\ }
\newcommand{\Subsection}{\subsection}
\providecommand{\nobreakdash}{\nobreak\hbox{-}}
\setlength{\emergencystretch}{2em}
\multlinegap0pt
\usepackage{amssymb}
\usepackage{float}
\usepackage{xcolor}
\newtheorem{lem}{Lemma}
\title{\textbf{Spectral moments and characteristic polynomials of vertex expansion hypergraphs of graphs}}
\author{Ge Lin, Changjiang Bu}
\date{Source: arXiv:2609.05897v1 (CC BY 4.0)}
\begin{document}
\maketitle

\noindent\emph{Source and licence.} \textbf{Spectral moments and characteristic polynomials of vertex expansion hypergraphs of graphs}. Version arXiv:2609.05897v1 (CC BY 4.0), distributed by arXiv under the Creative Commons Attribution 4.0 International licence, http://creativecommons.org/licenses/by/4.0/, as recorded in the arXiv licence metadata for this exact version and shown on the abstract page https://arxiv.org/abs/2609.05897v1. Full licence notice: \texttt{licenses/arxiv-2609.05897-CC-BY-4.0.txt}.\par\smallskip

\noindent\textit{Editorial scope.} Counted unit: the article's lem yinli4 (source0.tex lines 376--384). The article's complete original proof of it is delivered with it (source0.tex lines 386--440, one \texttt{proof} environment of 55 lines). No mathematics is written, repaired or re-derived: the statement and the proof are the article's own lines, in the article's own notation, and no retained line is substituted or re-flowed.  Retained uncounted as the source-local context the proof needs: lines 280--293.  Nothing was omitted from any retained span, so the package carries no omission record.  No retained line contains a citation, so the wrapper carries no bibliography and no bibliography entry.  No retained line contains a figure environment, an image-inclusion command or drawing code, so no image is delivered and the package declares no asset dependency.  Editorial formatting only, in addition: an AMS \texttt{amsart} preamble that restates the article's own declarations and macros used by the retained lines, namely the environments \texttt{lem} on separate counters, together with the standard packages those lines need.  The retained environments print in this document as Lemma 1 in order of appearance, whereas the article prints the same items with the numbering of its own full document.  The complete native source of the article as distributed in the arXiv e-print for this exact version is bundled unchanged as \texttt{sources/209473/source0.tex}.
\begin{lem}\label{yinli3}
Let $s \geq 2$ and $d \geq 2s$ such that $s \mid d$.
For $\widehat{G} = (\widehat{V},\widehat{E}) \in \mathcal{G}_{d}$,
let $D \in \mathfrak{D}_{d}(\widehat{G}^{[s]})$ and $f \in \mathcal{F}_{d}(\widehat{G}^{[s]})$ such that $D_{f} \cong D$.
For each $i \in \widehat{V}$ and any two distinct vertices $v,v' \in \mathbf{V}_{i}$,
we have
\begin{enumerate}

\item[(a)] $m_{D_f}(v,v')=m_{D_f}(v',v)$.

\item[(b)] $sm_{D_f}(v,v')=\sum_{\{i,j\} \in \widehat{E}}\sum_{u \in \mathbf{V}_{j}}m_{D_f}(u,v)=\sum_{\{i,j\} \in \widehat{E}}\sum_{u \in \mathbf{V}_{j}}m_{D_f}(u,v')$.

\end{enumerate}
\end{lem}

\begin{lem}\label{yinli4}
Let $s \geq 2$ and $d \geq 2s$ such that $s \mid d$.
For $\widehat{G} = (\widehat{V},\widehat{E}) \in \mathcal{G}_{d}$,
we have
\begin{align*}
\mathfrak{D}_{d}(\widehat{G}^{[s]})=
\bigcup_{Q \in \mathbf{Q}_{d,s}(\widehat{G})}\left\{ D : \mbox{$D_{f} \cong D$ for some $f \in \mathcal{F}_{d}(\widehat{G}^{[s]},Q)$} \right\}.
\end{align*}
\end{lem}

\begin{proof}
We prove this result by establishing the two inclusions.

For any $D \in \mathfrak{D}_{d}(\widehat{G}^{[s]})$,
let $f \in \mathcal{F}_{d}(\widehat{G}^{[s]})$ such that $D_{f} \cong D$.
We will show that there exists $Q \in \mathbf{Q}_{d,s}(\widehat{G})$ such that $f \in \mathcal{F}_{d}(\widehat{G}^{[s]},Q)$.
For each $i \in \widehat{V}$ and any two distinct vertices $v,v' \in \mathbf{V}_{i}$,
since $v'$ occurs as a non-root vertex in every rooted hyperedge of $f$ rooted at $v$,
the number of rooted hyperedges of $f$ rooted at $v$ is equal to $m_{D_f}(v,v')$.
It is known from Lemma \ref{yinli3}(a) that $m_{D_f}(v,v')=m_{D_f}(v',v)$.
Then the numbers of rooted hyperedges in $f$ rooted at each vertex in $\mathbf{V}_{i}$ are all equal,
and we denote it by $\mathrm{q}_{i}$.
Since $D_{f} \cong D$ is Eulerian,
we have $\mathrm{q}_{i} >0$.
Suppose that $q_{ij}$ (resp. $q_{ji}$) is the number of rooted hyperedges $\mathbf{V}_{i} \cup \mathbf{V}_{j}$ in $f$ rooted at vertices in $\mathbf{V}_{i}$ (resp. $\mathbf{V}_{j}$).
Since $H_{f} \cong \widehat{G}^{[s]}$,
we have $q_{ij}+q_{ji}>0$ for each $\{i,j\} \in \widehat{E}$ and $q_{ij}=q_{ji}=0$ for each $\{i,j\} \notin \widehat{E}$.
For each $i \in \widehat{V}$,
note that $\sum_{\{i,j\} \in \widehat{E}}q_{ij}=sm_{D_f}(v,v')=s\mathrm{q}_{i}$ and $\sum_{\{i,j\} \in \widehat{E}}q_{ji}=\sum_{\{i,j\} \in \widehat{E}}\sum_{u \in \mathbf{V}_{j}}m_{D_f}(u,v)=\sum_{\{i,j\} \in \widehat{E}}\sum_{u \in \mathbf{V}_{j}}m_{D_f}(u,v')$.
By Lemma \ref{yinli3}(b),
we have $\sum_{\{i,j\}\in\widehat{E}}q_{ji}=sm_{D_f}(v,v')$,
which yields that $$\sum_{\{i,j\} \in \widehat{E}}q_{ij}=\sum_{\{i,j\}\in\widehat{E}}q_{ji}=s\mathrm{q}_{i}.$$
%Denote $Q=(q_{ij})$ and $\mathbf{q}=(\mathrm{q}_{i})$.
Let $Q$ be the $|\widehat{V}| \times |\widehat{V}|$ matrix with entries $q_{ij}$,
and let $\mathbf{q}$ be the $|\widehat{V}|$-dimensional column vector with components $\mathrm{q}_{i}$.
Then we get $Q\mathbf{e}=Q^{\top}\mathbf{e}=s\mathbf{q}$.
The total multiplicity of arcs in $D_f$ is $$|E(D_f)|=\sum_{i \in \widehat{V}}s\mathrm{q}_{i}(2s-1)=(2s^{2}-s)\sum_{i \in \widehat{V}}\mathrm{q}_{i}.$$
Since $|E(D_f)|=d(2s-1)$,
we get $\sum_{i \in \widehat{V}}\mathrm{q}_{i}=\frac{d}{s}$,
that is, $\mathbf{e}^{\top}\mathbf{q}=\frac{1}{s}\mathbf{e}^{\top}Q\mathbf{e}=\frac{d}{s}$.
Thus,
we have $Q \in \mathbf{Q}_{d,s}(\widehat{G})$ and $f \in \mathcal{F}_{d}(\widehat{G}^{[s]},Q)$.
It implies that $$\mathfrak{D}_{d}(\widehat{G}^{[s]})\subseteq
\bigcup_{Q \in \mathbf{Q}_{d,s}(\widehat{G})}\left\{ D : \mbox{$D_{f} \cong D$ for some $f \in \mathcal{F}_{d}(\widehat{G}^{[s]},Q)$} \right\}.$$

For any $Q \in \mathbf{Q}_{d,s}(\widehat{G})$ and $f \in \mathcal{F}_{d}(\widehat{G}^{[s]},Q)$,
we next show that $D_{f}$ is Eulerian,
which implies that $D_{f} \cong D \in \mathfrak{D}_{d}(\widehat{G}^{[s]})$.
For each $i \in \widehat{V}$ and each $v \in \mathbf{V}_{i}$,
the number of rooted hyperedges in $f$ rooted at $v$ is $\frac{1}{s}(Q\mathbf{e})_{i}$.
Then we have $$\mathrm{deg}_{D_{f}}^{+}(v)=\frac{2s-1}{s}(Q\mathbf{e})_{i}.$$
On the other hand,
$v$ occurs as a non-root vertex in the rooted hyperedges of $f$ rooted at vertices in $\mathbf{V}_{i} \setminus \{v\}$,
with a total number of $\frac{s-1}{s}(Q\mathbf{e})_{i}$,
and also in the rooted hyperedges of $f$ rooted at vertices in $\mathbf{V}_{j}$ for $\{i,j\} \in \widehat{E}$,
with a total number of $q_{ji}$.
Then we have $$\mathrm{deg}_{D_{f}}^{-}(v)=\frac{s-1}{s}(Q\mathbf{e})_{i}+\sum_{\{i,j\} \in \widehat{E}}q_{ji}
=\frac{s-1}{s}(Q\mathbf{e})_{i}+\left(Q^{\top}\mathbf{e}\right)_{i}
=\frac{2s-1}{s}(Q\mathbf{e})_{i}.$$
Thus,
we get $\mathrm{deg}_{D_{f}}^{+}(v)=\mathrm{deg}_{D_{f}}^{-}(v)$,
that is, $D_{f}$ is Eulerian.
It implies that $$\mathfrak{D}_{d}(\widehat{G}^{[s]})\supseteq
\bigcup_{Q \in \mathbf{Q}_{d,s}(\widehat{G})}\left\{ D : \mbox{$D_{f} \cong D$ for some $f \in \mathcal{F}_{d}(\widehat{G}^{[s]},Q)$} \right\}.$$
\end{proof}

\end{document}
